
\subsubsection{2.1 RLHF Overoptimization and Reward Hacking}

Reinforcement learning from human feedback was introduced as a paradigm for aligning language models with human preferences through a reward model trained on preference annotations [Ouyang et al., 2022; Christiano et al., 2017]. RLHF has become the dominant alignment approach, enabling models that score high on helpfulness and harmlessness benchmarks [Bai et al., 2022; Stiennon et al., 2020]. However, the reward model is an imperfect proxy for human preference, and sustained optimization against this proxy causes the policy to exploit proxy weaknesses --- a phenomenon known as reward hacking or specification gaming [Krakovna et al., 2020; Skalse et al., 2022].

Coste et al. [2023] (arXiv:2310.02743) provide a clear empirical demonstration: as KL budget from the base policy increases from 0 to 8 nats, the reward model score rises monotonically while gold human preference peaks and then reverses. Gao et al. [2023] (arXiv:2210.10760, ICML 2023) characterize the scaling laws of reward model overoptimization at multiple RM sizes, showing consistent patterns across model scales. Both papers describe the divergence qualitatively --- RM rises while gold reverses --- but stop short of defining the \textit{normalized divergence gap} (RM\_norm $-$ gold\_preference) as a regression target or reporting cross-dataset slope comparison. The present work takes this step: OLS regression is fit to the normalized gap in both datasets, yielding $\beta$ = 0.1433 $nat^{-1}$ (Coste) and $\beta$ = 0.1599 $nat^{-1}$ (Gao), enabling the first cross-dataset characterization of the calibration-alignment divergence curve slope.

Ziegler et al. [2019] and Stiennon et al. [2020] document early cases where RLHF reward models diverge from human preferences in text summarization. These works establish the qualitative phenomenon; the contribution here is quantification and cross-dataset effect size comparison.

\subsubsection{2.2 Bidirectional Human-AI Alignment Measurement}

The ICLR 2025 Workshop on Bidirectional Human-AI Alignment synthesized 400 papers and found a systematic measurement asymmetry: virtually all major alignment benchmarks --- TruthfulQA \citep{Lin2022TruthfulQA}, BBQ \citep{Parrish2022BBQ}, HHH-RLHF, WinoBias \citep{Zhao2018Gender}, HELM \citep{Liang2023Holistic} --- measure the $AI\rightarrow Human$ direction exclusively. $Human\rightarrow AI$ alignment is studied in a parallel HCI and cognitive science track but has not been integrated with ML evaluation frameworks.

Lai et al. [2021] demonstrate that AI-assisted decision-making produces over-reliance: in several experimental settings, 60--80\% of humans agree with AI predictions regardless of AI accuracy. Bansal et al. [2021] and Buccinca et al. [2021] further characterize over-reliance patterns and propose interventions. However, these behavioral studies do not connect to RLHF optimization pressure --- they measure human calibration at a fixed model deployment snapshot, not as a function of how aggressively the model was optimized.

This work bridges these two tracks: RLHF overoptimization is treated as the empirical instantiation of bidirectional alignment tension. The scope is ''evaluation calibration divergence'' --- distinct from user behavioral calibration (appropriate reliance, Lai et al. [2021] paradigm). The connection to user over-reliance remains a theoretical extension for future behavioral work.

\subsubsection{2.3 Goodhart's Law and Proxy-Target Decoupling in ML}

Goodhart's Law --- ''when a measure becomes a target, it ceases to be a good measure'' --- has been formalized in the ML setting by Manheim and Garrabrant [2019] and others. Krakovna et al. [2020] compile a specification gaming list documenting empirical cases in RL. Skalse et al. [2022] formalize reward tampering and misspecification conditions under which proxy-target decoupling is guaranteed.

These theoretical works predict proxy-gold divergence under optimization pressure. The present work provides an empirical instantiation with specific quantitative parameters: in the RLHF setting with explicit reward model training, the proxy-gold gap grows at $\beta \approx$ 0.143--0.160 $nat^{-1}$ of KL budget. This connects the Goodhart/specification-gaming literature to a concrete, measurable phenomenon in modern RLHF training.

The three quantitative contributions --- the calibration-alignment divergence construct, the regression slope $\beta$ with cross-dataset replication, and the normalized gap metric --- are not derivable from any prior paper. The closest works are Coste et al. [2023] and Gao et al. [2023], which provide the raw experimental data; the contributions here are the regression framework, normalization methodology, cross-dataset comparison, and bidirectional alignment framing.

